\section*{Chapter \ref{ch:g10:functions} --- Functions}

\begin{solution}{exo:g10:functions:1}
$f(0) = 1$; \quad $f(2) = 8 - 6 + 1 = 3$; \quad
$f(-1) = 2 + 3 + 1 = 6$; \quad
$f\!\left(\tfrac12\right) = 2 \times \tfrac14 - \tfrac32 + 1 = \tfrac12 -
\tfrac32 + 1 = 0$.
\end{solution}

\begin{solution}{exo:g10:functions:2}
$f$: no constraint, \omterm{def:g10:functions:function}{domain} $\R$.

$g$: requires $x + 4 \neq 0$, \omterm{def:g10:functions:function}{domain} all reals except $-4$.

$h$: requires $2x - 6 \geq 0$, i.e.\ $x \geq 3$: \omterm{def:g10:functions:function}{domain}
$\intco{3}{+\infty}$.

$k$: $x^2 + 1 \geq 1 > 0$ never vanishes, \omterm{def:g10:functions:function}{domain} $\R$.
\end{solution}

\begin{solution}{exo:g10:functions:3}
\emph{\omterm{def:g10:functions:function}{Preimages} of $0$:} $x^2 - 2x = x(x-2) = 0$, so $0$ and $2$.

\emph{\omterm{def:g10:functions:function}{Preimages} of $3$:} $x^2 - 2x = 3$, i.e.\ $x^2 - 2x - 3 = 0$, i.e.\
$(x-3)(x+1) = 0$ (check by \omterm{def:g10:algebra:expand}{expanding}), so $3$ and $-1$.

\emph{\omterm{def:g10:functions:function}{Preimages} of $-1$:} $x^2 - 2x + 1 = (x - 1)^2 = 0$, so the single
\omterm{def:g10:functions:function}{preimage} $1$.
\end{solution}

\begin{solution}{exo:g10:functions:4}
$f(-1) = 1 + 2 + 1 = 4$: yes, $A(-1, 4)$ is on the \omterm{def:g10:functions:graph}{graph}.
$f(2) = 4 - 4 + 1 = 1$: yes, $B(2, 1)$ is on the \omterm{def:g10:functions:graph}{graph} too.
\end{solution}

\begin{solution}{exo:g10:functions:5}
\emph{1.} The largest value in the table is $4 = g(3)$ (\omterm{def:g10:functions:extrema}{maximum}), the
smallest is $-2 = g(0)$ (\omterm{def:g10:functions:extrema}{minimum}).

\emph{2.} On $\intcc{-3}{0}$ the \omterm{def:g10:functions:function}{function} decreases, and
$-1 < -0.5$, so $g(-1) \geq g(-0.5)$.

\emph{3.} On each \omterm{def:g10:numbers:interval}{interval} of monotonicity, $g$ takes each value at most
once, so $g(x) = 0$ has at most one solution per \omterm{def:g10:numbers:interval}{interval}: at most $3$
solutions in total. (Here $0$ lies between $-2$ and $1$, between $-2$ and
$4$, and between $0$ and $4$, so there are exactly three.)
\end{solution}

\begin{solution}{exo:g10:functions:6}
Let $u < v$. Then
\[
f(v) - f(u) = (-2v + 7) - (-2u + 7) = -2(v - u) < 0,
\]
since $v - u > 0$. So $f(v) < f(u)$: $f$ is strictly \omterm{def:g10:functions:variations}{decreasing} on $\R$.
\end{solution}

\begin{solution}{exo:g10:functions:7}
Complete the square: $(x+3)^2 = x^2 + 6x + 9$, so
\[
f(x) = x^2 + 6x + 11 = (x + 3)^2 + 2 .
\]
For every $x$, $(x+3)^2 \geq 0$, so $f(x) \geq 2$; and
$f(-3) = 0 + 2 = 2$. The \omterm{def:g10:functions:extrema}{minimum} of $f$ is $2$, attained at $x = -3$.
\end{solution}

\begin{solution}{exo:g10:functions:8}
\emph{1.} Length $+$ width $= 20$, so the length is $20 - x$ and
\[
A(x) = x(20 - x).
\]
Both dimensions must be positive: $0 < x < 20$.

\emph{2.} $A(4) = 64$, $A(8) = 96$, $A(10) = 100$, $A(12) = 96$,
$A(16) = 64$. The values rise until $x = 10$ then fall symmetrically:
the area seems largest for $x = 10$, i.e.\ for a \emph{square} garden.
(\cref{ch:g11:quad} proves this conjecture.)
\end{solution}

\begin{solution}{exo:g10:functions:9}
\emph{1.} For every $x$: $x^2 + 1 \geq 1 > 0$, so $f(x) = \frac{1}{x^2+1}$
is positive, and since $x^2 + 1 \geq 1$, taking inverses (both sides
positive) gives $f(x) \leq 1$.

\emph{2.} \omterm{def:g10:functions:extrema}{Maximum}: $f(0) = 1$ and $f(x) \leq 1$ for all $x$, so $1$ is the
\omterm{def:g10:functions:extrema}{maximum}, attained at $0$. \omterm{def:g10:functions:extrema}{Minimum}: $f(x) > 0$ always, but no value of $x$
achieves $0$ (a quotient of nonzero numbers is nonzero), and $f$ takes
values as close to $0$ as we like for large $x$; so $f$ has no \omterm{def:g10:functions:extrema}{minimum}.
\end{solution}

\begin{solution}{exo:g10:functions:10}
\emph{1.} $f(v) - f(u) = v^2 - u^2 = (v-u)(v+u)$. If $0 \leq u < v$, then
$v - u > 0$ and $v + u > 0$ (sum of a nonnegative and a positive number),
so $f(v) - f(u) > 0$: $f$ is strictly \omterm{def:g10:functions:variations}{increasing} on $\intco{0}{+\infty}$.

\emph{2.} If $u < v \leq 0$, then $v - u > 0$ still, but now
$v + u < 0$ (sum of a nonpositive and a negative number), so
$f(v) - f(u) < 0$: $f$ is strictly \omterm{def:g10:functions:variations}{decreasing} on $\intoc{-\infty}{0}$.
\end{solution}

\begin{solution}{pb:g10:functions:1}
\textbf{1.} \omterm{def:g10:functions:function}{Domain}: $x \neq 0$ (division by $x$). $f(1) =
\frac12(1 + 2) = \frac32$; $f(2) = \frac12(2 + 1) = \frac32$;
$f\!\left(\frac32\right) = \frac12\left(\frac32 +
\frac43\right) = \frac{17}{12}$.

\textbf{2.} $\frac12\left(x + \frac2x\right) = \frac32$ gives
$x^2 + 2 = 3x$, i.e.\ $x^2 - 3x + 2 = (x - 1)(x - 2) = 0$: the
\omterm{def:g10:functions:function}{preimages} of $\frac32$ are $1$ and $2$ --- as question 1
announced.

\textbf{3.} $f(x) = x$ gives $x + \frac2x = 2x$, so
$\frac2x = x$, $x^2 = 2$: \omterm{pb:g10:functions:1}{fixed points} $\sqrt2$ and $-\sqrt2$.
The positive one is the diagonal of the unit square, \omterm{ex:g10:numbers:classify}{irrational}
by the irrationality weekend problem of the Middle School volume.

\textbf{4.} $1 \to \frac32 \to \frac{17}{12} \to
\frac{577}{408}$: Heron's \omterm{def:g10:numbers:approx}{approximations} of $\sqrt2$. His
recipe ``average the guess with $2/$guess'' is exactly
\emph{iterating the \omterm{def:g10:functions:function}{function} $f$} --- and the number the
process chases is $f$'s \omterm{pb:g10:functions:1}{fixed point}.

\textbf{5.} The \omterm{pb:g10:functions:1}{fixed points} are the crossings of the \omterm{def:g10:functions:graph}{graph}
with the diagonal $y = x$: feeding outputs back in walks the
point along the \omterm{def:g10:functions:graph}{graph} towards that crossing. In
the two-thermometers weekend problem of the Middle School volume, the reading $-40^\circ$ was the same
crossing-the-diagonal idea for the conversion \omterm{def:g10:functions:function}{function}.

\textbf{6.} $f(v) - f(u) = \frac{v - u}{2} + \frac1v - \frac1u
= \frac{v - u}{2} + \frac{u - v}{uv}
= (v - u)\left(\frac12 - \frac{1}{uv}\right)
= \frac{(v - u)(uv - 2)}{2uv}$. For $0 < u < v \leq \sqrt2$:
$uv < 2$, the bracket is negative, $f(v) < f(u)$: \omterm{def:g10:functions:variations}{decreasing}.
For $\sqrt2 \leq u < v$: $uv > 2$: \omterm{def:g10:functions:variations}{increasing}. Turning point:
$\sqrt2$.

\textbf{7.} \omterm{def:g10:functions:variations}{Decreasing} before $\sqrt2$, \omterm{def:g10:functions:variations}{increasing} after: $f$
attains its \omterm{def:g10:functions:extrema}{minimum} at $\sqrt2$, worth
$f(\sqrt2) = \frac12\left(\sqrt2 + \frac{2}{\sqrt2}\right)
= \sqrt2$. Consequence: every output of the machine (positive
input) is at least $\sqrt2$ --- after one turn of the crank,
guesses live in $\intco{\sqrt2}{+\infty}$.

\textbf{8.} If $x > \sqrt2$, then $x^2 > 2$, so
$\frac2x < x$ and the average $f(x)$ of $x$ and $\frac2x$ is
less than $x$; and $f(x) \geq \sqrt2$ by question 7, with
equality only at $x = \sqrt2$. So
$\sqrt2 < f(x) < x$: each iteration strictly improves without
overshooting --- the sequence of question 4 slides down onto
$\sqrt2$.

\textbf{9.} On $\intoo{0}{+\infty}$: \omterm{def:g10:functions:variations}{decreasing} from
$+\infty$ (near $0$) down to the \omterm{def:g10:functions:extrema}{minimum} $\sqrt2$ at
$x = \sqrt2$, then \omterm{def:g10:functions:variations}{increasing} without bound.

\textbf{10.} $g(x) = x$ gives $x^2 = x$, $x(x - 1) = 0$: \omterm{pb:g10:functions:1}{fixed
points} $0$ and $1$. From $0.9$: $0.81$, $0.656$, $0.430$,
$0.185$ --- sliding towards $0$. From $1.1$: $1.21$, $1.464$,
$2.144$, $4.595$ --- fleeing to infinity. The \omterm{pb:g10:functions:1}{fixed point} $0$
attracts, $1$ repels: a hair's difference at the start, opposite
destinies.

\textbf{11.} $7 \to 22 \to 11 \to 34 \to 17 \to 52 \to 26 \to
13 \to 40 \to 20 \to 10 \to 5 \to 16 \to 8 \to 4 \to 2 \to 1$:
sixteen steps, peaking at $52$.

\textbf{12.} $27 \to 82 \to 41 \to 124 \to 62 \to 31 \to 94 \to
47 \to 142 \to 71 \to 214 \to \dots$ --- still climbing after
ten steps, on its way to a peak of $9\,232$ and a $111$-step
flight. Lesson: a two-line rule can produce behavior no one can
predict at a glance --- next-door numbers ($26$, $27$, $28$)
have wildly different flights.

\textbf{13.} A power of $2$ halves down its whole ladder:
$2^k \to 2^{k-1} \to \dots \to 2 \to 1$. At the bottom,
$h(1) = 4$, $h(4) = 2$, $h(2) = 1$: the flight enters the cycle
$4 \to 2 \to 1 \to 4$. \omterm{pb:g10:functions:1}{Fixed points} would need $\frac n2 = n$
(only $n = 0$) or $3n + 1 = n$ (negative): none among the
positive \omterm{def:g10:numbers:sets}{integers} --- the observed flights end not at a \omterm{pb:g10:functions:1}{fixed
point} but in that little \emph{cycle}.

\textbf{14.} A counterexample would be either a flight that
grows forever (never returning to $1$) or a second cycle
disjoint from $4 \to 2 \to 1$. Verification up to $10^{20}$
leaves infinitely many untested numbers --- exactly the
circle-regions lesson of the matchstick-oracle weekend problem of the Middle School volume: agreement on
finitely many cases proves nothing about all of them.

\textbf{15.} For Heron's machine we had \emph{structure} --- a
factored difference (question 6) revealing variations, a
\omterm{def:g10:functions:extrema}{minimum}, a squeeze --- while the hailstone's parity-switching
rule offers no such handle: nobody has found the structure that
tames it.

\textbf{16.} $\sqrt{x - 3}$: needs $x - 3 \geq 0$:
$\intco{3}{+\infty}$. $\frac{1}{x^2 - 9}$: needs
$x \neq \pm 3$: $\R$ deprived of $-3$ and $3$.
$\sqrt{9 - x^2}$: needs $x^2 \leq 9$: $\intcc{-3}{3}$.

\textbf{17.} $f(x) = 3$ gives $x + \frac2x = 6$, so
$x^2 - 6x + 2 = 0$, i.e.\ $(x - 3)^2 = 7$:
$x = 3 - \sqrt7$ or $x = 3 + \sqrt7$ (both positive: two exact
\omterm{def:g10:functions:function}{preimages}).

\textbf{18.} Physical sense while the stone is aloft:
$H(t) = 5t(4 - t) \geq 0$ for $t \in \intcc{0}{4}$. Completing:
$H(t) = 20 - 5(t - 2)^2$: \omterm{def:g10:functions:extrema}{maximum} $20$ m at $t = 2$ s.
Variations: \omterm{def:g10:functions:variations}{increasing} on $\intcc{0}{2}$ from $0$ to $20$,
\omterm{def:g10:functions:variations}{decreasing} on $\intcc{2}{4}$ back to $0$.

\textbf{19.} $20t - 5t^2 \geq 15 \iff t^2 - 4t + 3 \leq 0 \iff
(t - 1)(t - 3) \leq 0 \iff t \in \intcc{1}{3}$: two full
seconds above $15$ m.

\textbf{20.} Each machine has a \omterm{def:g10:functions:function}{domain} (where the rule is
allowed), a rule (the formula or the parity switch), and a
\omterm{def:g10:functions:graph}{graph} or table displaying it whole. Of each we asked the two
great questions: \emph{where does repetition lead} (\omterm{pb:g10:functions:1}{fixed
points}, cycles, attraction --- solved for Heron, open for the
hailstone) and \emph{how does the output move with the input}
(variations, extrema --- the stone's summit, Heron's \omterm{def:g10:functions:extrema}{minimum}).
The chapters ahead sharpen both: reference \omterm{def:g10:functions:function}{functions} next,
then, in later years, derivatives to measure variation and
limits to certify where iterations land.
\end{solution}
